% Chapter 5 content — included by main.tex; compile chapter5.tex for preview.

\section{Discussion}
\label{sec:chapter5}

\subsection{Applicability to Different Multi-Beam LEO Constellations}
\label{sec:discussion_applicability}

The proposed EABR framework is designed for general multi-beam LEO satellite systems and is not limited to a specific constellation. Although the main service-performance evaluation in this work uses a OneWeb-like constellation, the core idea of EABR depends on predictable satellite motion, beam footprints, beam-level EPFD contributions, and coverage overlap among neighboring LEO satellites. These properties are common in many multi-beam LEO systems.

Therefore, as long as the satellite geometry and beam coverage can be estimated in advance, the time-slot beam role record and candidate beam relation graphs can be constructed for different constellation configurations. The proposed framework can then identify the critical satellites and helper satellites according to the current system state. In this sense, EABR is applicable to different multi-beam LEO constellations.

\subsection{Helper Availability and Resource-Limited Recovery}
\label{sec:discussion_helper_resource}

Helper-based service recovery first requires neighboring satellites whose beams overlap with the closed beams of critical satellites. The availability of these helper satellite candidates depends on the constellation geometry, including the satellite distribution, orbital spacing, and beam footprint overlap. Section VI further examines helper satellite availability under different satellite distribution densities, and the results indicate that helper satellite candidates are generally available under the evaluated density scenarios.

Even when helper satellite candidates are available, the recovery performance of EABR is still limited by the available helper beam capacity and helper satellite transmit power. Therefore, insufficient recovery resources should be regarded as a system-level resource limitation rather than a failure of the reassociation framework. As long as feasible helper coverage and available power exist, EABR can use these resources for service recovery while maintaining the EPFD and power constraints.

\subsection{Priority-Aware Recovery and Starvation Risk}
\label{sec:discussion_priority_recovery}

In the HBR recovery score, user priority weights represent QoS-based service priorities. Under high-load conditions, the available helper recovery capacity may be insufficient to serve all closed-beam users. In this case, among the candidate closed-beam user subsets evaluated for each closed-recovery edge, subsets containing more high-priority users are more likely to obtain higher recovery scores and be selected first. Consequently, some low-priority users may remain unrecovered for several consecutive time slots, resulting in potential user starvation. This behavior reflects a tradeoff between priority-aware service recovery and user fairness. Future work may introduce a starvation-prevention mechanism to mitigate or avoid low-priority user starvation.